\documentclass[a4paper,12pt]{ctexart}
\usepackage[numbers, sort&compress]{natbib}
\usepackage{tocloft}
\usepackage{graphicx}
\usepackage{abstract}
\usepackage{geometry}
\usepackage{array}
\usepackage{datetime2}
\usepackage{fancyhdr}
\usepackage{ulem}  % 提供 \uline
\geometry{left=2.5cm, right=2.5cm, top=3cm, bottom=3cm} % 设定页边距
\CTEXsetup[format={\Large\bfseries}]{section}
%\setCJKmainfont{Source Han Serif SC}[BoldFont=Source Han Sans SC,ItalicFont=KaiTi] % 字体设置
\setCJKmainfont{SimSun}[BoldFont=SimHei,ItalicFont=KaiTi]
\ctexset{abstractname = {}} % 取消默认摘要标题
\ctexset{contentsname = {}} % 取消默认目录标题
\renewcommand{\cftsecleader}{\cftdotfill{\cftdotsep}}
\newcommand{\fixuline}[2][6cm]{\uline{\makebox[#1][c]{#2}}} % 自定义命令：固定长度下划线并居中显示内容
\newcommand{\upcite}[1]{\textsuperscript{\cite{#1}}} % 自定义引用上标命令


% !编译配方：xelatex*2
% !请务必选择该配方，若无该配方，请手动编译两次xelatex或自行定义该编译配方


% !声明论文主要信息
\newcommand{\LessonName}{你的课程名称}
\newcommand{\PaperType}{你的论文类型}
\newcommand{\MainTitle}{这里是大标题}
\newcommand{\SubTitle}{——这里是副标题}
\newcommand{\CollegeName}{这里是你的学院}
\newcommand{\AuthorName}{君の名前}
\newcommand{\ID}{1145141919810}
\newcommand{\Date}{\today} % 如需固定写死，可替换为yyyy-mm-dd
\newcommand{\Head}{这里是你的页眉}


\begin{document}

% --------封面--------
% 插入 logo，宽度设置为 \textwidth
\thispagestyle{empty}
\begin{center}
	\includegraphics[width=\textwidth]{logo.pdf}
\end{center}

\vspace{1cm}

% 课程名称
\begin{center}
	{\Huge \LessonName}\\
	\vspace{0.5cm}
	{\Huge \PaperType}
\end{center}

\vfill

% 主标题
\begin{center}
	{\huge \textbf{\MainTitle}}
\end{center}

% 副标题(不需要请自行注释掉)
\begin{center}
	{\LARGE \SubTitle}
\end{center}

\vfill

% 信息表格，放入 minipage 控制总宽度，避免偏移
\begin{center}
	\begin{minipage}{10cm}
		\renewcommand{\arraystretch}{1.8}
		\begin{tabular}{@{}>{\raggedleft}p{2.5cm}@{\hspace{-0.2cm}}>{\centering\arraybackslash}p{6.7cm}@{}}
			学院： & \fixuline{\CollegeName} \\
			姓名： & \fixuline{\AuthorName} \\
			学号： & \fixuline{\ID} \\
		\end{tabular}
	\end{minipage}
\end{center}

\vfill

% 日期
\begin{center}
	{\large \Date}
\end{center}

% --------摘要--------
\newpage
%\pagenumbering{Roman} % 开始罗马数字页码
\begin{abstract}
    \thispagestyle{empty}
    \begin{center}
        {\bfseries\Large 摘\hspace{2em}要} \\[1em]
    \end{center}
	这里是摘要

	\vspace{2em}

    {\noindent\bfseries 关键词：} \textit{关键词1；关键词2}
\end{abstract}

% --------目录--------
\newpage
\begin{center}
	{\bfseries\Large 目\hspace{2em}录}
	\tableofcontents % 生成目录
	\thispagestyle{empty}
\end{center}

% --------引言--------
\newpage
\cleardoublepage
\fancypagestyle{myheader}{
	\fancyhf{}                                  % 清除默认设置
	\fancyhead[C]{\Head} % 页眉左侧内容（可自定义）
	\fancyhead[R]{第\thepage 页}                % 页眉右侧页码（可自定义）
}
\pagestyle{myheader}  % 启用 fancy 页眉
\renewcommand{\headrulewidth}{0.4pt} % 页眉横线宽度
\pagenumbering{arabic} % 设置页码格式是阿拉伯数字
\setcounter{page}{1}


\section*{引言} % 星号表示不编号
\addcontentsline{toc}{section}{引言} % 手动添加目录项
这里是引言


% --------正文--------
\section{第一节}
这里是第一节\upcite{m1}


\section{第二节}
这里是\textbf{第二节}\upcite{m1,m2}


\section{第三节}
这里是\textit{第三节}\upcite{m1,x1,m2}

% --------结语--------
\section*{结语} % 带空格的标题
\addcontentsline{toc}{section}{结语} % 目录项
这里是结语

% --------参考文献--------
\renewcommand{\refname}{参考文献}
\addcontentsline{toc}{section}{参考文献}
\begin{thebibliography}{99}  
	\bibitem{m1} 参考文献1
	\bibitem{x1} 参考文献2
	\bibitem{m2} 参考文献3
\end{thebibliography}

\end{document}